\documentclass[11pt]{article}
\usepackage[T1]{fontenc}
\usepackage{textcomp}
\usepackage[margin=0.75in]{geometry}
\usepackage{amsmath,amssymb,amsthm}
\usepackage{array,booktabs}
\usepackage{hyperref}
\setlength{\parindent}{0pt}
\newtheorem{theorem}{Theorem}
\theoremstyle{definition}
\newtheorem{definition}{Definition}
\title{Flow Proof Artifact: prop-22-triangle-from-three-lines}
\author{Generated by \texttt{flow doc proof}}
\date{Flow Proof Book}
\begin{document}
\maketitle
To construct a triangle from three straight lines, each less than the sum of the other two.\\[0.5em]
\textbf{Source.} euclid — Elements, Book I, Proposition 22\\[0.5em]
\bigskip
\noindent{\large \textbf{Derived fact 1.} \textit{To construct a triangle from three straight lines, each less than the sum of the other two} \hfill the Euclidean plane \cdot Euclid Book I \cdot Proposition 22: construct a triangle from three lines}\par
\smallskip\noindent\textit{Coordinate: the Euclidean plane \cdot Euclid Book I \cdot Proposition 22: construct a triangle from three lines}\par
\smallskip\noindent\textit{Source: euclid — Elements, Book I, Proposition 22}\par
\smallskip\noindent\textit{Built on: proposition 1: equilateral triangle on a segment, for Book I of the Elements in on the Euclidean plane, proposition 20: the sum of two sides exceeds the third, for Book I of the Elements in on the Euclidean plane}\par
\medskip
\noindent\textbf{Goal.} To construct a triangle from three straight lines, each less than the sum of the other two\par
\noindent$\displaystyle \text{triangle from lines LMN exists}$\par
\medskip
\noindent
\begin{tabular}{@{} >{\raggedright\arraybackslash}p{0.06\textwidth} >{\raggedright\arraybackslash}p{0.47\textwidth} >{\raggedleft\arraybackslash}p{0.41\textwidth} @{}}
\textbf{\#} & \textbf{Proof} & \textbf{Mathematics} \\
\midrule
\textbf{1.} & We prove that proposition 22: construct a triangle from three lines for Book I of the Elements in the Euclidean plane. & \\[0.45em]
\textbf{2.} & Let circle\_center\_A\_radius\_F = circle\_at\_A\_with\_radius\_F. & \\[0.45em]
\textbf{3.} & Let circle\_center\_B\_radius\_G = circle\_at\_B\_with\_radius\_G. & \\[0.45em]
\textbf{4.} & Let K = intersection\_of\_circles\_at\_A\_and\_B. & \\[0.45em]
\textbf{5.} & We invoke the derived fact: lines\_F\_G\_H\_satisfy\_triangle\_inequality. & \\[0.45em]
\textbf{6.} & We invoke the axiom governing Book I of the Elements in the Euclidean plane: proposition 1: equilateral triangle on a segment, for Book I of the Elements in on the Euclidean plane. & \\[0.45em]
\textbf{7.} & We invoke the derived fact governing Book I of the Elements in the Euclidean plane: proposition 20: the sum of two sides exceeds the third, for Book I of the Elements in on the Euclidean plane. & \\[0.45em]
\textbf{8.} & From step 2, step 3, step 4, step 5, step 6, and step 7, by the chain of results established in Book I (step 2, step 3, step 4, step 5, step 6, and step 7), we obtain triangle from lines LMN exists. Hence proven. & $\displaystyle \text{triangle from lines LMN exists}$ \\[0.45em]
\bottomrule
\end{tabular}
\label{thm:Geometry-Euclid Book I-proposition_22_construct_a_triangle_from_three_lines}
\par\medskip\hrule
\end{document}
